\section{Motivation}
\begin{frame}{Overview}{}
    \begin{block}{Joint Work -- Its not me alone}
        \centering
        Thorsten Kranz, Gregor Leander, Ko Stoffelen, Friedrich Wiemer
        \includegraphics[width=0.35\textwidth]{data/logo_rub}
        \hspace{2em}
        \includegraphics[width=0.45\textwidth]{data/ru_en_branding}
    \end{block}
    \begin{block}{Outline}
        \vspace{0.5em}
%        \begin{itemize}
%            \item Preliminaries
%            \item State of the Art
%            \item Related Work
%            \item Updated State of the Art
%        \end{itemize}
        \tableofcontents
    \end{block}
\end{frame}

\begin{frame}{What is the \textsc{xor} count,}{and why is it important?}
    \begin{block}{Some facts}
    \begin{itemize}
        \visible<1->{\item Lightweight Block Ciphers}
        \visible<1->{\item Efficient Linear Layers}
        \visible<1->{\item MDS matrices are \enquote{optimal} (regarding security)\footnote{Are they?}}
    \end{itemize}
    \begin{itemize}
        \visible<2->{\item What is the lightest implementable MDS matrix?}
        \visible<2->{\item What about additional features (Involutory)?}
    \end{itemize}
    \end{block}
\end{frame}

\section{Preliminaries}
\begin{frame}{What is an MDS matrix?}{}
%    \begin{columns}
%        \begin{column}{0.45\textwidth}
    \begin{block}{Definition: MDS}
        A matrix $M$ of dimension $k$ over the field $\mathbb{F}$ is \emph{maximum distance separable} (MDS),
        iff all possible submatrices of $M$ are invertible (or nonsingular).
    \end{block}
%        \end{column}
%        \begin{column}{0.45\textwidth}
    \visible<2->{\begin{example}
        The \textsc{Aes MixColumn} matrix is defined over $\mathbb{F}_{2^8} \cong \mathbb{F}[x]/\mathtt{0x11b}$:
        \begin{equation*}
            \begin{pmatrix}
                \mathtt{0x02} & \mathtt{0x03} & \mathtt{0x01} & \mathtt{0x01} \\
                \mathtt{0x01} & \mathtt{0x02} & \mathtt{0x03} & \mathtt{0x01} \\
                \mathtt{0x01} & \mathtt{0x01} & \mathtt{0x02} & \mathtt{0x03} \\
                \mathtt{0x03} & \mathtt{0x01} & \mathtt{0x01} & \mathtt{0x02}
            \end{pmatrix}
            =
            \begin{pmatrix}
                x & x + 1 & 1 & 1 \\
                1 & x & x + 1 & 1 \\
                1 & 1 & x & x + 1 \\
                x + 1 & 1 & 1 & x
            \end{pmatrix}
        \end{equation*}
        This is a (right) \emph{circulant} matrix: $\text{circ}(x, x+1, 1, 1)$.
    \end{example}}
%        \end{column}
%    \end{columns}
\end{frame}
\note{\begin{itemize}
        \item The degree of $\mathbb{F}$ normally fits to the S-box width.
        \item The dimension $k$ resembles the state.
\end{itemize}}

\begin{frame}{What is an MDS matrix?}{Constructions}
    \centering
    \begin{tikzpicture}[%
        mindmap,
        level 1 concept/.append style={%
            level distance=130,
            sibling angle=45,
        },
        scale=0.7, every node/.style={scale=0.7},
    ]

    % Kind of attacks
    \begin{scope}[%
        mindmap,
        concept color=saphierblau,
        text=white,
    ]
        \node [concept] {\textbf{\large MDS}}[clockwise from=45]
        child  {node [concept] (construction) {\textbf{by construction}}[clockwise from=20]
            child {node [concept] (cauchy) {\textbf{Cauchy}}}
            child {node [concept] (vandermonde) {\textbf{Vander\-monde}}}
        }
        child [visible on=<2->] {node [concept] (other) at (2.5, -1) {\textbf{other}}[clockwise from=30]
            child [visible on=<2->] {node [concept] (circ) {\textbf{Circulant}}}
            child [visible on=<2->] {node [concept] (hadamard) {\textbf{Hadamard}}}
            child [visible on=<2->] {node [concept] (toeplitz) {\textbf{Toeplitz}}}
        }
        child [visible on=<3->] {node [concept] (subfield) {\textbf{Subfield}}};
    \end{scope}

    \begin{pgfonlayer}{background}
        \draw (vandermonde)+(2,0) node[scale=0.6] {can be made involutory} ;
        \draw [visible on=<2->] (circ)+(2,0) node[scale=0.6] {used in \textsc{Aes}, \textsc{Gr\o{}stl}, \textsc{Whirlpool}, \textsc{Qarma}, \textsc{Midori}} ;
        \draw [visible on=<2->] (hadamard)+(2,0) node[scale=0.6] {can be made involutory, used in \textsc{Anubis}, \textsc{Khazad}, \textsc{Joltik}} ;
        \draw [visible on=<3->] (subfield)+(2,0) node[scale=0.6] {used in \textsc{Whirlwind}} ;
    \end{pgfonlayer}

    \end{tikzpicture}
\end{frame}

\begin{frame}{What is an MDS matrix?}{Representations}
\end{frame}

\begin{frame}{What is an MDS matrix?}{Representations}
\end{frame}

\begin{frame}{The General Linear Group}{}
\end{frame}

\section{State of the Art and Related Work}
\begin{frame}{State of the Art}{Before our Paper}
\end{frame}

\begin{frame}{Related Work I}{}
\end{frame}

\begin{frame}{Related Work II}{}
\end{frame}

\begin{frame}{State of the Art}{After our Paper}
\end{frame}

\section{Future Work}
\begin{frame}{Future Work}{}
\end{frame}
